\section{離散対称性と CPT}

\noindent\textbf{この章の中心命題}
\[
\boxed{
\begin{gathered}
C:\text{粒子}\leftrightarrow\text{反粒子},
\qquad
P:\mathbf{x}\mapsto-\mathbf{x},
\qquad
T:t\mapsto -t,\\
CPT\ \text{は局所相対論的量子場理論の対称性}
\end{gathered}
}
\]
相対論的場の量子論では、電荷反転 $C$、空間反転 $P$、時間反転 $T$ を個別に調べる必要がある。自然界の相互作用はそれらを必ずしも別々には保たないが、通常の局所相対論的量子場理論では組み合わせた $CPT$ は保たれる。

これまで扱ってきた対称性の多くは連続的であった。空間の平行移動、回転、ローレンツ変換、位相変換は、連続的なパラメータを少しずつ変えられる。これに対して、空間を丸ごと反転する、粒子を反粒子に入れ替える、時間の向きを反転する、といった変換は連続的には単位変換へ戻せない。このような対称性を離散対称性という。

\subsection{空間反転 $P$}
空間反転 $P$ は、時刻をそのままにして空間座標の符号を反転する変換である。
\[
P:\quad (t,\mathbf{x})\mapsto(t,-\mathbf{x}).
\]
スカラー場なら
\[
\phi(t,\mathbf{x})\mapsto \phi(t,-\mathbf{x})
\]
と変換する。ベクトル場では、普通の空間ベクトルは符号を反転するが、角運動量や磁場のような軸性ベクトルは符号の変わり方が異なる。この違いは、弱い相互作用を考えるときに重要になる。

ディラック場では、空間反転はスピノール成分にも作用し、代表的には
\[
\psi(t,\mathbf{x})
\mapsto
\gamma^0\psi(t,-\mathbf{x})
\]
と書ける。ここで $\gamma^0$ が現れるのは、ディラック方程式の形を保つためである。

\begin{experimentalfact}{パリティ非保存}
\factitem{操作}{低温で偏極させたコバルト60原子核のベータ崩壊を観測し、原子核スピンの向きと放出電子の方向の相関を測定する。}
\factitem{前提}{空間反転が弱い相互作用の対称性なら、スピン方向に対して電子が出やすい向きと出にくい向きは反転後に区別できない。}
\factitem{結果}{電子は原子核スピンの向きに対して非対称に放出される。}
\factitem{示唆}{弱い相互作用は空間反転 $P$ を対称性として保たない。}
\end{experimentalfact}

この実験事実は、「物理法則は鏡に映しても同じである」という素朴な期待が、弱い相互作用では成り立たないことを示した。

\subsection{電荷反転 $C$}
電荷反転 $C$ は、粒子と反粒子を入れ替える変換である。電磁場では、電荷の符号を反転すると電場・磁場の源も反転するので、
\[
A_\mu\mapsto -A_\mu
\]
と変換する。ディラック場では、粒子を消す演算子と反粒子を作る演算子が入れ替わる。記法としては、電荷共役場
\[
\psi^c:=C\bar{\psi}^{\,T}
\]
を用いて表す。ここで右辺の $C$ は電荷共役行列であり、添字をもつスピノール成分を反粒子のスピノールへ対応させるための行列である。

QED では、電子と陽電子の電荷が反対であることを同時に反映させれば、電荷反転の下で理論の形は保たれる。一方、弱い相互作用では左巻き粒子と右巻き粒子の扱いが異なるため、$C$ も一般には対称性ではない。

\subsection{時間反転 $T$}
時間反転 $T$ は
\[
T:\quad (t,\mathbf{x})\mapsto(-t,\mathbf{x})
\]
である。古典力学では、速度や運動量の符号を反転すれば、時間反転された運動を考えられる。量子論ではさらに、時間発展の位相
\[
e^{-iEt}
\]
を正しく反転させる必要があるため、$T$ は単なる線形演算子ではなく反ユニタリ演算子として表される。ここで反ユニタリとは、複素数 $i$ を $-i$ に移す作用を含むことを意味する。

この点は形式的に見えるが、物理的には重要である。シュレーディンガー方程式
\[
i\frac{\partial}{\partial t}\psi(t)=\hat{H}\psi(t)
\]
で $t\mapsto -t$ としたとき、左辺の $i$ も同時に複素共役で符号を変えなければ、方程式の形が保たれないからである。

\subsection{CP 対称性とその破れ}
$C$ と $P$ がそれぞれ破れても、組み合わせた $CP$ は保たれるのではないか、という可能性がある。しかし弱い相互作用では $CP$ も厳密な対称性ではない。

\begin{experimentalfact}{中性 K 中間子の CP 破れ}
\factitem{操作}{中性 K 中間子の崩壊を観測し、長寿命成分がどの終状態へ崩壊するかを測定する。}
\factitem{前提}{CP が厳密な対称性なら、長寿命成分 $K_L$ は CP の性質により二つのパイ中間子へは崩壊しない。}
\factitem{結果}{$K_L$ が二つのパイ中間子へ崩壊する事象が観測される。}
\factitem{示唆}{弱い相互作用は $CP$ も厳密には保たない。}
\end{experimentalfact}

この事実は、粒子と反粒子の振る舞いが完全には対称でないことを意味する。宇宙に物質が反物質より多く残っている理由を考える上でも、$CP$ 破れは重要な役割をもつ。

\subsection{CPT 定理}
個々の $C$、$P$、$T$、あるいは組み合わせた $CP$ が破れても、通常の局所相対論的量子場理論では、三つをすべて組み合わせた $CPT$ は対称性として保たれる。この主張を CPT 定理という。

ここで「通常の局所相対論的量子場理論」とは、少なくとも次の構造をもつ理論を指す。
\[
\text{ローレンツ不変性},
\qquad
\text{局所性},
\qquad
\text{ユニタリな時間発展},
\qquad
\text{安定な真空}.
\]
この条件のもとで、粒子と反粒子は同じ質量をもち、寿命も対応する過程で等しくなる。たとえば電子と陽電子の質量が等しいことは、CPT 対称性と整合している。

本ノートでは CPT 定理の証明には踏み込まない。重要なのは、CPT が経験的に「たまたまよく成り立っている」対称性ではなく、局所性とローレンツ不変性をもつ量子場理論の深い帰結として現れる点である。したがって、もし CPT の破れが観測されれば、単に一つの相互作用の性質が分かったというだけでなく、局所相対論的量子場理論の前提そのものを見直す必要がある。
